%************************************************
\chapter{Introduction}\label{cap:intro}
%************************************************

% switch between these depending on page size, or modify directly
% same for the rest of the chapters
\tikz[remember picture,overlay] \node[opacity=1,inner sep=0pt] at (current page.center){\includegraphics[width=\paperwidth,height=\paperheight]{./Figures/cover/port_title.png}};
%\tikz[remember picture,overlay] \node[opacity=0.3,inner sep=0pt] at ([yshift=3cm,xshift=2cm]current page.center){\includegraphics[width=\paperwidth,height=\paperheight]{./Figures/cover/sesshu_1.jpg}};
\clearpage 

% PREFACE? espacio en blanco?

\vspace*{100pt}
\label{sec:preface}
\lettrine{I}{t was the year of} \textsc{1868} when the first traffic signals in
 the world came to be, installed outside the Houses of Parliament in London.
 The manually operated gas-lit signals that the railway engineer J.P. Knight
 designed were not, however, intended to control an intersection between busy
 flows of horse-drawn carriages, but rather to allow the safe crossing of a
 street: the reason to be of the first traffic signal was to serve the pedestrian.
 % mover ref fuera d aki a stat of the art o algo

Along with the evolution and expansion of the automobile, traffic signals also
 evolved and became more ubiquitous in cities around the world. The shiny novel
 vehicles quickly became traffic, a fast and dangerous new kind of it, and that
 traffic needed to be controlled; traffic signals —or traffic lights, or
 \textit{robots}, as known in South Africa— were the logical choice for that.\par

The high regard for the pedestrian that gave the traffic light its conception
 was quickly forgotten in the process. As streets and roads surrendered more
 and more space to the car, more and more engineers and architects preached its
 new rule as a symbol of progress. For many, the priority was to move cars, not
 people, and that reflected on traffic signals, as will be seen in the pages
 that follow.

Despite the shift that mobility has experienced ever since, placing other modes
 of transport in their rightful place, some of those habits and designs
 continue to be present in our cities and traffic signals are no exception.

\medskip

\noindent{When I moved to València for the start of my civil engineering
 education, traffic signals quickly became an unexpected topic in which to
 ponder. While in my native Palma I was used to waiting for the green light
 before crossing the street, to do so in València felt absurd. When
 walking in Palma, a red light meant that cars were in the way: waiting for the
 green one was not a matter of compliance but of self-preservation; in València,
 a red light as a pedestrian meant that \textit{maybe} some cars could appear,
 almost always in a rhythm and quantity that did not call for the length of the
 halt.}

It is well known that traffic control devices influence people's
 behaviour,\footnote{
    For the case of traffic signals, see the \hyperref[sec:motivation]{following section} and the \hyperref[subsec:bckg_sota]{review of the state of the art} for examples.
 } and unsurprisingly, I frequently saw valencians jaywalking to avoid long
 waits in front of an empty roadway. This practice was not limited to people
 in a hurry or to rebellious types: families with children, elderly people,
 parents with strollers and in general anyone who otherwise would be expected
 to be a law-abiding\footnote{
    Pursuant to article 145 of the Spanish traffic law
     (\textit{Reglamento General de Circulación}~\cite{ReglamentoGeneraldeCirculacion}),
     ignoring a red light as a pedestrian is a «severe» infraction, punishable
     with a fine of 200 euros. In València, I never saw the police prosecuting
     anyone for it, even when the jaywalking was done directly in front of the
     officers.}
 citizen would usually just check if traffic was clear and
 cross, ignoring the red light. The refusal to follow the traffic signals 
 extended beyond pedestrians, as cyclists and motorized scooter users often did
 the same. I observed that, while more widespread in side streets and other lightly trafficked ways, these attitudes were still common even in wide avenues where the crossing involved many lanes of traffic at high speeds; thoroughfares which in other cities would have been perceived as
 dangerous to cross without clearance.
 
It felt logical: if during everyday walks or bike rides people constantly
 encountered traffic signals ordering them to stop, sometimes reasonably, 
 sometimes for periods double the time that the motorized traffic required,
 sometimes at places where the actual necessity of the traffic signal itself
 came into question; a red light would lose its meaning to them, and hence,
 its use as a means of protecting people.

I unfortunately saw many accidents resulting from pedestrians and cyclists not
 following traffic signals during my studies in València, most of the time just
 a fright that left everyone unharmed, but occasionally the consequences were much more severe.
 
This situation felt not only unfair but dangerous: the burden of being cautious rested \textit{both} on those ignoring the red light and those where actually given the right of way, the drivers. Any lack of visibility or attention from both sides could result in the loss of life, even when one of them should, according to the system, not expect any conflicting traffic.

\medskip
\noindent When, as part of the European Mobility Week, I paid a visit in 2023 to the traffic management centre of the city,\footnote{More information about this centre is available in the official website of the city council~\cite{aytoval}.} I quickly realized that this was not a result of design limitations or a lack of traffic data; on the contrary, it was a deliberate choice. 

Personnel in the centre explained  that an adaptive traffic signal control system switched traffic signal phases in real time using the traffic counts of more than 3000 induction loop detectors in the car lanes, so —as they added with a hint of pride— the system was able to adapt to the traffic conditions of the city in real time.

Having counts of motorized traffic as the only input does not necessarily imply that the system will disregard other modes completely, however, when I asked about how pedestrians and cyclists were considered and if they were aware of the amount of jaywalking happening because of signal timings, the tone in the answer changed.

The technicians assured that the system also accounted for other modes, giving examples of some avenues having green waves were optimized for cyclist speeds —although still using car counts as input— and some busy city centre streets with longer pedestrian phases during shopping hours. Interestingly, the smaller network of bike counting loop detectors was not mentioned at all. They also claimed that pedestrians and cyclists not respecting red lights were a problem of education and enforcement that happened in any city, and it was not a task for the traffic signal control system to address.

As if the place of modes different than automobile in the traffic signal control system was not sufficently clear, the big screens overshadowing the control room made it manifest. On the sides, a rotation of traffic camera feeds displayed intersections and tunnels, placed on tall masts so vehicular approaches could be seen, but turning pedestrians and cyclists as blurry objects appearing from under the treetops; in the middle, the street network was displayed as a cobweb of coloured lines using the data that the induction loops measured for automobiles: València reduced into a diagram of car flows.


The visit continued, but in the process, more and more citizens expressed with their questions and comments how they felt that the showcased system did not serve them correctly as users of the streets, which seemed certainly contrasting with the satisfaction that the personnel expressed about it. The claims regarding how the system optimized timings in real time clashed with several complaints regarding waiting times for drivers and non-drivers alike, and while the personnel had boasted of the city being the third city in the world by number of traffic signals,\footnote{A pride that they have shared with the press, see~\cite{cadenaser}.} some visitors viewed the amount of signals themselves excessive.

While the personnel brushed off these concerns, an increasing undertone of dissatisfaction lingered over the visit, and I could not help but remember it every time I found myself waiting in a red light that felt unnecessarily long.

\medskip

\noindent This visit has remained a cautionary tale in my mind ever since, one that I have often recalled during my civil engineering education.

Engineering is about solving problems, and engineers develop complex systems to do so, but it is easy to become so focused on their abstractions in our pursuit of solutions that we lose sight of the problems that needed to be solved in the first place.

Infrastructure is meant to serve people, and not the opposite. The traffic control system that I had the opportunity to see seemed like an excellent solution to the problem of optimizing vehicular flow, but did not look like a solution for how citizens in general moved about the city.

When professor Rodrigues presented the topic of this thesis, I could not help but remember how I felt bearing the long red lights of València and how they continued to serve only cars, just because pedestrian and cyclist flows still lack their proper place in mainstream traffic signal control systems.

The real-world aspect of it enticed me even further, as I felt that facing a real-world scenario —the city of Copenhagen— would surely bring complexity and challenges, but also opportunities in which to find engineering solutions that actually explore beyond ideal abstractions and serve people.

\medskip

\noindent While this thesis does not aim to revolutionize traffic signal control, it does aim to be a useful contribution towards making traffic engineering better for all users of the streets.

\medskip

% decir q nos metemos en nuestras sims y nuestros datos t nos ovlidaos de usuario

% decir q esto lo he recalleao incluso haciendo este tfm

%blabla q luego con la motorizacion nos olvidamos de peatones y bicis y tal pero
% que en inicio la cosa era multimodal

\section{Motivation} \label{sec:motivation}

As the preface explains, traffic lights influence their users' behaviour, and the way they are designed and operated can have a significant impact on the safety, efficiency, and overall experience of people.

Beyond personal experiences, there are several reasons that motivate the work presented in this thesis.

\medskip

\noindent Increasing urbanization and population growth have led to the growth of cities worldwide, and with them, the need for capable transportation systems. Relying solely on the automobile for moving people in an urban environment is not an efficient solution, and it also results in negative externalities such as congestion, pollution, and accidents. Consequently, cities are progressively promoting alternative modes of transport, such as public transport, cycling, and walking, to reduce the reliance on private vehicles.

However, to turn these modes a fully viable alternative, infrastructure must be developed accordingly. It is not only safety and comfort standards that matter, but especially the travel efficiency also plays a key role in attracting daily commuters to use a specific mode \cite{Rietveld2004,Brunner2024}.

Hence, as cities become more diverse in terms of transportation modes, the need for effective \acrfull{tsc} systems that can manage efficiently and safely this multimodal traffic becomes more and more important. 

Walking and cycling offer sustainable and promising alternatives to the automobile in terms of efficiency, especially when considering the constraints of an urban environment: a 3\,metre-wide car lane in a city can move around 1600 persons per hour; a bike lane of the same width can move around 3800 people, and an equivalent sidewalk, as much as 9000~\cite{natco2013}. In addition of their health and environmental benefits, these modes offer probably the best cost-efficiency ratio in terms of space and money per person transported.

\medskip

\noindent If a \gls{tsc} system is to be effective, it must account for the behaviour of all road users, including pedestrians and cyclists. This is not reduced solely to the safety of these vulnerable modes, however. As will be seen further in Section~\ref{subsec:bckg_sota}, pedestrian and cyclist dynamics have a certain impact in the design and performance of traffic signals, and with it in the overall transport system in general \cite{Noland1996}.

\medskip

\noindent Regarding moving by foot, pedestrian crossing speed influences the determination of several \gls{tsc} parameters, such as pedestrian clearance time or cycle length. However, unlike some traffic manuals suggest, it is not a fixed value, but rather a variable that depends on the characteristics of the pedestrians themselves. Pedestrian crossing speed is affected by various factors, including age, health condition, physical fitness, and psychological expectations; hence significant differences in crossing speeds exist among different individuals. Relevant studies~\cite{Huang2010, Chandra2013} proved that pedestrian crossing speed at intersections typically follows a normal distribution.

Pedestrian dynamics are also affected by the presence of other pedestrians~\cite{Alhajyaseen2009}. When demand is high, the difference in speeds between different persons has a significant impact, as faster pedestrians are forced to slow down due to the presence of slower ones. The influence of other pedestrians is particularly relevant in the case of bi-directional flows, precisely the case of signalized crosswalks, where pedestrians from both sides of the street interact with each other. The interaction between opposing pedestrian flows can lead to congestion and delays, as pedestrians must navigate around each other while crossing the street. This can result in longer crossing times and increased waiting times for pedestrians, especially during peak hours when pedestrian demand is high.

On the other hand, as it was discussed already in the \hyperref[sec:preface]{preface}, jaywalking is of major importance for \gls{tsc} in particular and urban mobility in general, both in terms of efficiency as of safety. Research ~\cite{Crompton1979, Gaarder1989, Martin2006} has shown that most pedestrians are willing to wait around 20 to 30 seconds at a signalized intersection before compliance starts to drop, albeit this value is highly dependent on the context, with factors such as the presence of other pedestrians, the perceived safety of the crossing, or the traffic culture of the city \cite{Retzko1974}. Furthermore, pedestrians crossing streets in non-signalized crossings —be it in zebra crossings lacking signals, using in a mixed-traffic street or simply jaywalking— creates opportunities for interaction between modes, which even in non-risky situations imply changes of speed and behaviour for all users involved that cannot be ignored when modelling urban traffic \cite{Marisamynathan2018}.

\medskip

\noindent Similarly, taking cyclists into account is also relevant in the design of a proper \gls{tsc} system, especially in an urban setting.

Its ability to keep high flows even under congested conditions makes cycling a very attractive mode of transport for the urban environment \cite{Brunner2024}. However, the dynamics of cyclists are different from those of motorized vehicles, and they require different considerations in terms of traffic signal timing and phasing. 

Akin to pedestrians, cyclist dynamics are also very heterogeneous and influenced by factors such as age, gender, experience or bicycle type \cite{PérezCastro2025}. Cyclists also have different acceleration and deceleration characteristics than motorized vehicles, which can affect their ability to clear intersections quickly; and may show specific behaviours in manoeuvres such as left turns. Additionally, and also like pedestrians, cyclists are vulnerable road users in traffic conflicts with other modes —such as turning vehicles— which can impact their safety and travel time, hindering the perception of the bicycle as a modal choice \cite{Apostola2014}. Thus, these factors must be taken into account when designing street infrastructure, including \gls{tsc}.

Cyclists are especially sensitive to becoming stopped by red lights, as each stop implies a significant loss of speed and energy, which can be discouraging for them. This is particularly relevant in urban areas that require cyclists to navigate through multiple intersections, as frequent stops can significantly increase travel time and reduce the attractiveness of cycling as a mode of transport \cite{Apostola2014}. In contrast, a well-designed \gls{tsc} system has been proven to improve the efficiency of cycling as a mode and attract users towards the bicycle \cite{Aultman-Hall1997, Rietveld2004}.

\medskip

\noindent Despite these considerations, and as the \hyperref[subsec:bckg_sota]{review of the state of the art} shows further in this thesis, most of the research on \gls{tsc} has focused on motorized traffic, with limited attention given to pedestrians and cyclists.

When in the 1990s, the first \acrfull{rl} methods were applied to \gls{tsc}, the focus stayed on optimizing traffic flow for motorized vehicles, and further progress in the field of \acrfull{rltsc} has largely stayed centred only in the automobile.\footnote{Precisely, in a recent revision of the state of the art \cite{Xiao2024}, these challenges were explicitly identified as directions for future work.} On the other hand, most of the \gls{rltsc} research focuses in relatively simple scenarios, which may not accurately represent the complexity of real-world urban traffic networks, and this is specially true in the limited corpus of multimodal \gls{rltsc} literature. Furthermore, works treating the inclusion of bicycle traffic are very scarce in the current \gls{rltsc} state of the art.

\medskip

\noindent Considering the above, it is clear that there is a need for further research on \gls{tsc} that takes into account the dynamics of pedestrians and cyclists, especially in realistic urban scenarios. This thesis aims to address this gap.

\section{Problem statement}

While \acrfull{rl} has already been proven as a robust option for \acrfull{tsc}, most of the research performed so far has focused exclusively on automobiles, as the later \hyperref[sec:lit_review]{literature review} shows, and the work done considering multimodal traffic flows has mostly limited itself to simple intersections and not to realistic depictions of urban networks. Furthermore, among the studies accounting for multimodal traffic, the inclusion of cyclists still has not received the necessary attention if we account for their relevance in urban mobility.

\medskip

\noindent Considering this, the problem statement of this thesis can be summarized as follows:

\indent \textit{Development and evaluation of reinforcement learning methods for traffic signal control that account for pedestrians and cyclists in realistic scenarios.}

\section{Research questions}\label{sec:rqs}

The present thesis aims to provide address the stated problem by solving a number of research questions:

\begin{enumerate}[I.]
   \item Is it possible to develop a scenario for traffic control simulation that incorporates realistic demand and infrastructure based in a real-world setting?
   \item Can existing reinforcement learning methods for traffic signal control be adapted or extended to account for multimodal traffic flows, including pedestrians and cyclists?
   \item How do multimodal reinforcement learning methods for traffic signal control perform in realistic scenarios?
   \item Which influence have the architectural decisions taken in the design of multimodal reinforcement learning methods for traffic signal control on their performance across modes?
\end{enumerate}

\section{Scope and limitations}

While the topic of multimodal traffic signal control is broad and complex, the scope of this thesis is necessarily limited by our time and capabilities. As such, the following limitations are acknowledged:

\begin{itemize}
    \item The nature of the traffic signal control methods developed in this work is limited to reinforcement learning, and does not consider other traffic management strategies or technologies.
    \item The multimodal aspect of this thesis is centred on pedestrians and cyclists for the scenario design and the traffic signal control methods. While highly relevant for traffic signal control, other relevant modes such as public transport, autonomous vehicles, emergency vehicles, or micromobility devices are explicitly excluded from the scope of this thesis.
    \item Although aiming to be realistic, the development and evaluation of the traffic signal control methods is limited to simulation, and does not include real-world implementation or testing.
    \item The scenarios and methods developed in this thesis are based on the city of Copenhagen, Denmark, and may not be directly applicable to other cities or regions with different traffic patterns, infrastructure, or regulations.
\end{itemize}